% Figura extraída de "Tema 3.tex" (línea 220 del original)
\documentclass[12pt,border=1pt]{standalone}
\usepackage[spanish,es-tabla]{babel}
\usepackage[utf8]{inputenc}
\usepackage{amsmath}
\usepackage{xcolor}
\usepackage{tikz}
\usepackage{pgfplots}
\usepackage{circuitikz}
\usetikzlibrary{babel, circuits.ee}

\definecolor{applegreen}{rgb}{0.55, 0.71, 0.0}
\definecolor{airforceblue}{rgb}{0.36, 0.54, 0.66}
\definecolor{amber}{rgb}{1.0, 0.49, 0.0}

\newcommand{\diff}[3][1]{%
	\IfBlankTF{#1}{\frac{\mathrm{d}{#2}}{\mathrm{d}{#3}}}%
	{\frac{\mathrm{d}^{#1} {#2}}{\mathrm{d}{#3}^{#1}}}%
}%
\newcommand{\pdiff}[3][1]{%
	\IfBlankTF{#1}{\frac{\partial {#2}}{\partial {#3}}}%
	{\frac{\partial^{#1} {#2}}{\partial {#3}^{#1}}}%
}%

\begin{document}
\begin{tikzpicture}
\draw[ultra thick] (-3,2)--(-1.2,2); % Unión recta
     \draw[ultra thick] (-1.2,2) to[out=-20, in=120] (0,1); % Unión suavizada
    \draw[ultra thick] (0,1) to[out=-60, in=180] (1.7,0.5); % Desplazado 2.3 unidades hacia abajo
        \node at (-2,2.5) {P};
    \node at (2,2.5) {N};
    \draw[ultra thick] (1.7,0.5)--(3,0.5);
\begin{scope}[xscale=-1]
        \draw[ultra thick] (-3,1.8)--(-1.2,1.8); % Unión recta
        \draw[ultra thick] (-1.2,1.8) to[out=-20, in=120] (0,1); % Unión suavizada reflejada
        \draw[ultra thick] (0,1) to[out=-60, in=180] (1.7,0.2); % Desplazado 2.3 unidades hacia abajo y reflejado
        \draw[ultra thick] (1.7,0.2)--(3,0.2); % Reflejado
    \end{scope}    
    \draw[-latex, thick] (0,-0.2)--(0,3);
    \draw[-latex, thick] (-3,0)--(3,0);
        \draw[thick] (-1.2, 2.5)-- (-1.2, -0.2) node[below] {$x=x_\mathrm{p}$};
    \draw[thick] (1.2, 2.5)-- (1.2, -0.2)node[below] {$x=x_\mathrm{n}$};
    \end{tikzpicture}
\end{document}
